% =============================================================================
% exercises/ejercicios_cap06.tex
% Ejercicios propuestos — Capítulo 6: Contaminación Atmosférica
% =============================================================================

\section*{Ejercicios Propuestos}
\addcontentsline{toc}{section}{Ejercicios propuestos}
\label{sec:ejercicios_cap06}

\begin{bloque_ejercicios}[Ejercicios --- Cap\'{\i}tulo 5: Contaminaci\'on Atmosf\'erica]

\begin{ejercicios}

% ----- Ejercicio 1: Fuentes y clasificación de contaminantes ----------------
\item Clasifique cada uno de los siguientes contaminantes indicando:
(i)~su \textbf{origen} (biol\'ogico, tierra s\'olida, oc\'eano o formaci\'on
in situ), (ii)~si es \textbf{primario o secundario}, y
(iii)~si pertenece a los \textbf{contaminantes criterio} o no:

{%
\centering
\begin{tabular}{clccc}
\toprule
\textbf{N\'um.} & \textbf{Contaminante} & \textbf{Origen} &
\textbf{Primario/Secundario} & \textbf{Criterio} \\
\midrule
a) & \ce{SO2} (erupci\'on volc\'anica) & & & \\
b) & \ce{O3} troposf\'erico        & & & \\
c) & \ce{CH4} (humedales)          & & & \\
d) & \ce{NO} (combusti\'on)         & & & \\
e) & \ce{PAN}                       & & & \\
f) & \ce{H2SO4} (aerosol)           & & & \\
g) & \ce{Pb} (gasolina)             & & & \\
h) & DMS (fitoplancton)             & & & \\
\bottomrule
\end{tabular}
\par
}%

% ----- Ejercicio 2: Modelo de difusión de COS -------------------------------
\item Con base en el modelo de intercambio de COS presentado en el
cap\'{\i}tulo, en el que la concentraci\'on decrece exponencialmente con la
profundidad $z$ seg\'un $C(z)=C(0)\exp(-z/H)$, con altura de escala
$H=\left(D/k\right)^{1/2}$:

\begin{enumerate}[label=\alph*)]
    \item Si el coeficiente de difusi\'on turbulenta se duplica a
          $D = \SI{0.010}{\metre\squared\per\second}$ y la constante de
          destrucci\'on permanece en $k = 5.0\times10^{-6}\,\si{\per\second}$,
          calcule la nueva altura de escala $H'$ y la profundidad a la que
          $C(z) = C(0)/2$.

    \item Demuestre que la vida media de una mol\'ecula de COS en el oc\'eano
          es siempre $\tau = 1/k$, independientemente del valor de $D$.
          \textit{(Siga el procedimiento del ejemplo del cap\'{\i}tulo.)}

    \item ?`C\'omo afectar\'{\i}a una disminuci\'on de la productividad biol\'ogica
          marina al flujo de COS hacia la atm\'osfera? Argumente
          cualitativamente en t\'erminos de $C(0)$ y la densidad de columna.
\end{enumerate}

\textit{[Respuesta parcial: (a) $H' = \SI{44.7}{\metre}$;
profundidad para $C(0)/2$: $z = H'\ln 2 \approx \SI{31.0}{\metre}$]}

% ----- Ejercicio 3: Ciclo geoquímico del CO2 --------------------------------
\item El ciclo de meteorizaci\'on de la calcita (\ce{CaCO3}) involucra la
reacci\'on neta:
\[
    \ce{CaCO3 (s) + CO2 (g) + H2O (l) <=> Ca^{2+} (ac) + 2HCO3^- (ac)}
\]

\begin{enumerate}[label=\alph*)]
    \item Identifique en qu\'e direcci\'on procede la reacci\'on
          (derecha o izquierda) durante: (i) la meteorizaci\'on continental
          por lluvias \'acidas, y (ii)~la precipitaci\'on de sedimentos
          marinos. Justifique.

    \item Cuando el \ce{CO2} se disuelve en agua pura se forman tres
          especies en equilibrio. Escriba las tres ecuaciones de equilibrio
          en secuencia y la reacci\'on neta, tal como se presenta en el
          ejemplo del cap\'{\i}tulo.

    \item El tiempo de residencia del \ce{HCO3^-} en los oc\'eanos es
          $\tau \approx 7.5\times10^4$~a\~nos. Si la tasa de entrada de
          \ce{HCO3^-} al oc\'eano es $F = \SI{3.0e12}{\mole}/$~a\~no,
          estime la cantidad total de \ce{HCO3^-} en el oc\'eano usando
          $M = F \cdot \tau$.

    \item Discuta cu\'al ser\'{\i}a el efecto de un incremento sostenido de
          \ce{CO2} atmosf\'erico (debido a actividades antropog\'enicas)
          sobre el equilibrio de la reacci\'on neta anterior y sobre
          la acidez de los oc\'eanos.
\end{enumerate}

\textit{[Respuesta parcial:
(c) $M = 3.0\times10^{12} \times 7.5\times10^{4}
    = 2.25\times10^{17}\,\si{\mole}$]}

% ----- Ejercicio 4: Carboxihemoglobina (CO) ---------------------------------
\item La concentraci\'on de carboxihemoglobina (\%COHb) en sangre se
estima con la ecuaci\'on:
\[
    \%\ce{COHb} = 0.005\,[\ce{CO}]^{0.85}\,(\alpha t)^{0.63}
\]
donde $[\ce{CO}]$ est\'a en ppm, $\alpha$ es el coeficiente de actividad
y $t$ es el tiempo de exposici\'on en minutos.

\begin{enumerate}[label=\alph*)]
    \item Una persona sedentaria ($\alpha = 0.94$) es expuesta a
          $[\ce{CO}] = 35$~ppm durante 2 horas.
          Calcule \%COHb e indique qu\'e efectos fisiol\'ogicos
          se esperan seg\'un los umbrales del cap\'{\i}tulo.

    \item Un trabajador de construcci\'on ($\alpha = 3.0$) labora en un
          espacio confinado con $[\ce{CO}] = 50$~ppm durante 90~min.
          Calcule \%COHb y eval\'ue el riesgo.

    \item Usando los resultados de los incisos anteriores, explique por
          qu\'e el mismo nivel de \ce{CO} resulta m\'as peligroso para
          personas con actividad f\'{\i}sica intensa que para personas
          en reposo.

    \item Estacionamiento subterr\'aneo. La norma mexicana NOM-021 fija
          el l\'{\i}mite de \%COHb en trabajadores en $\leq 5\%$.
          ?`Cu\'al es la concentraci\'on m\'axima permisible de \ce{CO} (en ppm)
          para un trabajador sedentario ($\alpha = 0.94$) expuesto 8 horas?
          \textit{Despeje $[\ce{CO}]$ de la ecuaci\'on.}
\end{enumerate}

\textit{[Respuesta parcial:
(a) $\%\text{COHb} \approx 3.1\%$, efectos ligeros;
(b) $\%\text{COHb} \approx 7.5\%$, efectos psicomotores;
(d) $[\ce{CO}]_{\max} \approx 17$~ppm]}

% ----- Ejercicio 5: Compuestos de azufre y lluvia ácida ---------------------
\item El di\'oxido de azufre emitido a la atm\'osfera se oxida
progresivamente y puede generar lluvia \'acida.

\begin{enumerate}[label=\alph*)]
    \item Escriba la secuencia de reacciones por las que el \ce{SO2}
          atmosf\'erico se convierte en \'acido sulf\'urico (\ce{H2SO4})
          en presencia de agua y ox\'{\i}geno. Indique en cu\'al paso
          interviene el \ce{OH} radical.

    \item Una central termoelect\'ectrica quema carb\'on con contenido
          de azufre del 2.5\% en masa. Si la tasa de consumo es de
          \SI{500}{\mega\gram\per\hour}, estime la emisi\'on de \ce{SO2}
          en \si{\mega\gram\per\hour}. Dato: masa molar \ce{S} = 32~g/mol;
          \ce{SO2} = 64~g/mol.

    \item Seg\'un los umbrales de efectos en la salud del cap\'{\i}tulo,
          clasifique el riesgo para la poblaci\'on circundante si
          la concentraci\'on ambiente de \ce{SO2} alcanza 0.5~ppm
          durante 3 horas.

    \item Explique c\'omo el \ce{SO2} contribuye al incremento de
          PM$_{10}$ y PM$_{2.5}$, identificando la especie qu\'{\i}mica
          intermedia formada.
\end{enumerate}

\textit{[Respuesta parcial:
(b) Masa de S quemado = $500 \times 0.025 = \SI{12.5}{\mega\gram\per\hour}$;
emisi\'on de \ce{SO2} $= 12.5 \times (64/32) = \SI{25}{\mega\gram\per\hour}$;
(c) $0.5 > 0.3$~ppm $\Rightarrow$ riesgo alto: enfermedades
cardiorrespiratorias]}

% ----- Ejercicio 6: Ozono troposférico y estratosférico ---------------------
\item Con base en el \textbf{Cuadro~\ref{ozon:1}} y el texto del cap\'{\i}tulo:

\begin{enumerate}[label=\alph*)]
    \item Complete la siguiente tabla comparativa indicando
          \textit{altitud}, \textit{porcentaje del ozono total},
          \textit{efecto principal} y \textit{problem\'atica actual}
          para cada tipo de ozono:

{%
\centering
\begin{tabular}{lcc}
\toprule
\textbf{Caracter\'{\i}stica} & \textbf{Troposf\'erico} & \textbf{Estratosf\'erico}\\
\midrule
Altitud (km)         & & \\
\% del total         & & \\
Efecto principal     & & \\
Problem\'atica actual & & \\
\bottomrule
\end{tabular}
\par
}%

    \item La formaci\'on de ozono troposf\'erico requiere de COV,
          NO$_x$ y luz solar. Explique por qu\'e en una ciudad con
          alta emisi\'on de NO pero baja de COV la concentraci\'on de
          \ce{O3} puede ser menor que en los suburbios.

    \item Mencione dos factores meteorol\'ogicos u orogr\'aficos que
          favorecen la acumulaci\'on de ozono troposf\'erico en
          ciudades como la Ciudad de M\'exico o Los \'Angeles,
          con base en la informaci\'on del cap\'{\i}tulo.

    \item El ozono estratosf\'erico absorbe la radiaci\'on UV.
          ?`Qu\'e consecuencias para la salud humana se derivar\'{\i}an
          de una reducci\'on significativa de esta capa?
          Mencione al menos dos efectos.
\end{enumerate}

% ----- Ejercicio 7: Partículas y visibilidad --------------------------------
\item El alcance \'optico $V$ (visibilidad horizontal) se relaciona
con el coeficiente de extinci\'on $\beta_{ext}$ mediante la f\'ormula de
Koschmieder:
\[
    V = \frac{3.912}{\beta_{ext}} \quad [\si{\kilo\metre}]
\]

\begin{enumerate}[label=\alph*)]
    \item Calcule la visibilidad en un d\'{\i}a pristino con
          $\beta_{ext} = 10\ \text{Mm}^{-1}$ y en un episodio de
          contaminaci\'on severa con
          $\beta_{ext} = 1{,}000\ \text{Mm}^{-1}$.

    \item Seg\'un el cap\'{\i}tulo, la humedad relativa mayor a 70\%
          incrementa $\beta_{ext}$ entre 2 y 5 veces por crecimiento
          higrosc\'opico de las part\'{\i}culas. Si en una ma\~nana con
          HR~$= 85\%$ el $\beta_{ext}$ seco es de $200\ \text{Mm}^{-1}$,
          estime el intervalo de visibilidad esperado.

    \item Las PM$_{2.5}$ tienen di\'ametros de \SIrange{0.12}{0.22}{\micro\metre},
          que coinciden con la condici\'on $\alpha \approx 1$ para luz
          visible ($\lambda \approx \SI{0.55}{\micro\metre}$).
          Explique en t\'erminos del par\'ametro de Mie
          $\alpha = 2\pi r/\lambda$ por qu\'e estas part\'{\i}culas
          dispersan la luz de manera m\'axima.

    \item Clasifique los tres gases (\ce{SO2}, \ce{NO}, \ce{H2O})
          que menciona el cap\'{\i}tulo como influyentes en la visibilidad,
          describiendo el mecanismo espec\'{\i}fico de cada uno.
\end{enumerate}

\textit{[Respuesta parcial:
(a) $V_{pristino} = 3.912/10 \times 10^{-3}
    = \SI{391.2}{\kilo\metre}$;
    $V_{contam} = \SI{3.9}{\kilo\metre}$;
(b) $\beta_{ext}^{humedo} \in [400, 1000]\ \text{Mm}^{-1}$;
    $V \in [3.9, 9.8]\ \si{\kilo\metre}$]}

% ----- Ejercicio 8: Plomo en la atmósfera -----------------------------------
\item El tetraetilo de plomo (\ce{(C2H5)4Pb}) se a\~nad\'{\i}a a la gasolina
a raz\'on de \SI{1.1}{\gram}/gal (equivalente a\SI{0.29}{\gram\per\litre}).

\begin{enumerate}[label=\alph*)]
    \item Si un veh\'{\i}culo consume \SI{40}{\litre}/semana de gasolina
          con esta formulaci\'on, calcule la masa de plomo emitida
          a la atm\'osfera por semana. Dato: fracci\'on m\'asica de Pb en
          \ce{(C2H5)4Pb} $= 207/(207 + 4 \times 29) = 64.1\%$.

    \item Un ni\~no de 5 a\~nos presenta concentraci\'on sangu\'{\i}nea de
          plomo de \SI{15}{\micro\gram\per\deci\liter}. Seg\'un los umbrales
          del cap\'{\i}tulo (CDC 2023), ?`este valor est\'a por debajo o
          por encima del umbral de da\~nos hematol\'ogicos?
          ?`Qu\'e efectos neurol\'ogicos se podr\'{\i}an observar?

    \item Explique el mecanismo por el cual el pH del agua potable
          influye en la concentraci\'on de plomo ingerido a partir
          de tuber\'{\i}as de plomo. ?`Un agua con pH~$= 6$ es m\'as o
          menos riesgosa que una con pH~$= 8$?

    \item Desde el punto de vista de la pol\'{\i}tica ambiental, ?`por
          qu\'e la eliminaci\'on del plomo en la gasolina se considera
          una de las medidas de salud p\'ublica de mayor impacto en
          el siglo XX? Argumente con base en las v\'{\i}as de exposici\'on
          descritas en el cap\'{\i}tulo.
\end{enumerate}

\textit{[Respuesta parcial:
(a) Masa TEL = $40 \times 0.29 = \SI{11.6}{\gram}/$semana;
masa Pb $= 11.6 \times 0.641 = \SI{7.4}{\gram}/$semana;
(b) \SI{15}{\micro\gram\per\deci\liter} $<$ \SI{20}{\micro\gram\per\deci\liter}: por debajo de umbral
hematol\'ogico, pero en zona de riesgo neurol\'ogico]}

% ----- Ejercicio 9: Esmog y contaminantes primarios/secundarios -------------
\item El esmog fotoqu\'{\i}mico se genera cuando contaminantes primarios
reaccionan bajo la influencia de la luz solar.

\begin{enumerate}[label=\alph*)]
    \item Describa las diferencias fundamentales entre el
          \textbf{esmog industrial} (tipo Londres) y el
          \textbf{esmog fotoqu\'{\i}mico} (tipo Los \'Angeles) en cuanto a:
          condiciones meteorol\'ogicas, contaminantes involucrados
          y efectos en la salud.

    \item La reacci\'on simplificada de formaci\'on de esmog
          fotoqu\'{\i}mico puede escribirse como:
          \[
            \text{COV} + \text{NO}_x + h\nu \longrightarrow
            \text{esmog fotoqu\'{\i}mico} (\ce{O3},\,\text{PAN},\ldots)
          \]
          Indique qu\'e es el PAN\index{PAN}, a qu\'e clase de contaminante
          pertenece (primario o secundario) y por qu\'e resulta
          especialmente irritante.

    \item En episodios de esmog, las concentraciones de \ce{NO}
          son m\'aximas en la ma\~nana y decaen al mediod\'{\i}a,
          mientras que el \ce{O3} presenta el comportamiento opuesto.
          Explique este ciclo diurno en t\'erminos de las reacciones
          del cap\'{\i}tulo.

    \item ?`Por qu\'e las ciudades con cuencas rodeadas de monta\~nas
          (como la Ciudad de M\'exico) son especialmente vulnerables
          a los episodios de esmog fotoqu\'{\i}mico?
\end{enumerate}

% ----- Ejercicio 10 (avanzado): Compuestos tóxicos y análisis de riesgo ----
\item \textbf{[Avanzado]} El benceno (\ce{C6H6}) es uno de los compuestos
t\'oxicos listados en el \textbf{Cuadro~\ref{airtox}} y es un agente
canser\'ogeno reconocido.

\begin{enumerate}[label=\alph*)]
    \item Clasifique al benceno seg\'un las dos categor\'{\i}as de
          contaminantes descritas en la Secci\'on~\ref{ctox}
          (\textit{contaminante criterio} vs.\ \textit{t\'oxico ambiental})
          e indique el efecto espec\'{\i}fico sobre la salud reportado en
          el \textbf{Cuadro~\ref{efectos}}.

    \item El benceno proviene principalmente de la evaporaci\'on
          de gasolina y de la combusti\'on incompleta. Identifique
          en qu\'e secci\'on del cap\'{\i}tulo se menciona la quema de
          biomasa como fuente de \ce{CH3Cl} y relacione ese proceso
          con la emisi\'on de benceno: ?`ambos son fuentes
          biol\'ogicas, antropog\'enicas o mixtas?

    \item La metodolog\'{\i}a EPA~IRIS establece una concentraci\'on
          de riesgo unitario (URF) para benceno de
          $\text{URF} = 7.8\times10^{-6}\,(\mu\text{g/m}^3)^{-1}$.
          El riesgo individual de c\'ancer se calcula como:
          \[
              R = \text{URF} \times C \times \frac{ED}{70}
          \]
          donde $C$ es la concentraci\'on en $\mu$g/m$^3$,
          $ED$ es la duraci\'on de la exposici\'on en a\~nos y 70 es
          la expectativa de vida. Calcule $R$ para una persona que
          vive 30 a\~nos en una zona con $C = \SI{5}{\micro\gram\per\cubic\metre}$.

    \item Si la EPA considera ``aceptable'' un riesgo individual
          $R \leq 1\times10^{-6}$ (1 caso en un mill\'on),
          ?`el valor calculado en el inciso anterior es aceptable?
          ?`Qu\'e concentraci\'on m\'axima de benceno ser\'{\i}a permisible
          para cumplir este criterio en las mismas condiciones
          de exposici\'on?
\end{enumerate}

\textit{[Respuesta parcial:
(c) $R = 7.8\times10^{-6} \times 5 \times (30/70)
    \approx 1.67\times10^{-5}$;
(d) No aceptable; $C_{\max} = 10^{-6}/(7.8\times10^{-6}\times 30/70)
    \approx \SI{0.30}{\micro\gram\per\cubic\metre}$]}

\end{ejercicios}

\end{bloque_ejercicios}
% =============================================================================
% FIN DE ejercicios_cap06.tex
% =============================================================================
